\section{A sharp formal obstruction at weight six}
\label{depthnew:sec:obstruction}

The audited manuscript proves a complementary-depth bound for one-variable
multiple polylogarithms: after adjoining a transformed argument and endpoint
constants, a weight-$w$, depth-$d$ value has a representation of reduced
depth at most $w-d$. The first critical triple weight is therefore six.
The following result identifies exactly what survives if one tries to
reduce every such triple further using shuffle products and the six
standard argument transformations. It is a theorem about a specified
formal quotient, with no assertion of numerical period independence.

\subsection{The quotient and the argument action}

Let $V=\mathbb Qx+\mathbb Qy$, let $\operatorname{Sh}(V)$ be its shuffle
algebra, and let $J$ be the augmentation ideal. Set
\[
 Q_6=(J/J^2)_6.
\]
The depth of a binary word is its number of $y$ letters. Including words
with trailing $x$ makes linear substitutions immediate; shuffle
normalization with $H_x(z)=\log z$ removes such trailing letters without
changing the positive-weight indecomposable quotient.
The usual outer-letter-first interpretation is
\[
 x=\frac{dt}{t},\qquad y=\frac{dt}{1-t},\qquad
 H_{0^{s_1-1}1\cdots0^{s_d-1}1}(z)
 =\Li_{s_1,\ldots,s_d}(z,1,\ldots,1),
\]
where $0$ denotes $x$ and $1$ denotes $y$.

The group of six M\"obius transformations permuting $0,1,\infty$ acts
on the letter space through the matrices
\begin{equation}\label{depthnew:eq:matrices}
 R=\begin{pmatrix}0&-1\\-1&0\end{pmatrix},\qquad
 M=\begin{pmatrix}1&0\\1&-1\end{pmatrix}.
\end{equation}
Matrices act on columns, so $R(x)=-y$, $R(y)=-x$, whereas
$M(x)=x+y$, $M(y)=-y$. These are the pullbacks for $z\mapsto1-z$
and $z\mapsto z/(z-1)$, respectively. The six matrices are
$I,R,M,RM,MR,RMR$.
Landen substitution fixes the initial endpoint zero. Complementation
changes it to one; ordinary path composition supplies endpoint constants
and products of lower-weight integrals. Thus the same letter action
models these argument changes modulo products and endpoint constants.
The formal space is defined independently of any special-point evaluation.

\begin{theorem}[Exactly one missing triple direction]
\label{depthnew:thm:codimension}
Let $D_2\subset Q_6$ be the span of words of depth at most two, and
let $E_2$ be the span of their images under the six matrices above. Then
\[
 \dim Q_6=9,\qquad \dim E_2=8.
\]
The span is unchanged if all of $\operatorname{GL}_2(\mathbb Q)$ is
allowed in place of the six matrices. Its annihilator is generated by
coefficient pairing with the Lie polynomial
\begin{equation}\label{depthnew:eq:omega}
 \Omega=\bigl[[x,[x,y]],[y,[x,y]]\bigr].
\end{equation}
On the three triple generators used in the manuscript this pairing is
\begin{equation}\label{depthnew:eq:triple-values}
\begin{array}{c|ccc}
\text{word}&000111&001011&001101\\
\text{index}&(4,1,1)&(3,2,1)&(3,1,2)\\ \hline
\langle\Omega,\text{word}\rangle&0&-1&2.
\end{array}
\end{equation}
In particular, the class of $(3,2,1)$ supplies the single additional
formal direction.
\end{theorem}

\begin{proof}
The coefficient pairing is
$\langle\sum a_ww,\sum b_ww\rangle=\sum a_wb_w$.
Lie polynomials are primitive for the coproduct dual to shuffle
multiplication: this follows from
$\Delta x=x\otimes1+1\otimes x$, the same formula for $y$, and closure
of primitive elements under commutators. Hence every Lie polynomial
pairs to zero with every shuffle product of two nonempty words.
Consequently $\Omega$ defines a functional on $Q_6$.

Every word in $\Omega$ has exactly three $x$ and three $y$ letters,
so the functional vanishes on $D_2$. It is nonzero: the coefficient
of $001011$ is $-1$, and expansion gives the three entries in
\eqref{depthnew:eq:triple-values}.
There is also a covariance under every invertible linear substitution
$N$. Put $C=[x,y]$, $A=[x,C]$, and $B=[y,C]$.
The element $C$ transforms by $\det N$, while the pair $(A,B)$
transforms by $\det N$ times the original two-dimensional matrix.
Their commutator therefore satisfies
\begin{equation}\label{depthnew:eq:covariance}
 N(\Omega)=(\det N)^3\Omega.
\end{equation}
Since $N$ and its transpose have the same determinant,
\[
 \langle\Omega,N(w)\rangle
 =\langle N^{\mathsf T}(\Omega),w\rangle
 =(\det N)^3\langle\Omega,w\rangle.
\]
Thus $\Omega$ annihilates the full linear-letter orbit of $D_2$.
Both proposed spans have codimension at least one.

For the reverse inequality, the Lyndon description of the shuffle
algebra \cite{depthnew:Radford} gives the binary Witt count
\[
 \dim Q_6=\frac{2^6-2^3-2^2+2}{6}=9.
\]
The nine Lyndon words are
\[
\begin{gathered}
000001,\ 000011,\ 000101,\ 000111,\ 001011,\ 001101,\\
001111,\ 010111,\ 011111.
\end{gathered}
\]
For each Lyndon word $l$, let $b(l)$ denote its standard Lie bracketing:
if $l=uv$ and $v$ is its longest proper Lyndon suffix, put
$b(l)=[b(u),b(v)]$, with the letters as initial cases.
These nine primitive elements form a dual family of functionals on $Q_6$.

Take, in order, the eight word classes
\[
\begin{gathered}
 I(000001),\ I(000011),\ I(000101),\\
 R(000001),\ R(000011),\ R(000101),\\
 M(000001),\ M(000101),
\end{gathered}
\]
and pair them with the eight $b(l)$ obtained by omitting $001101$
from the ordered Lyndon list. Direct commutator expansion gives
\begin{equation}\label{depthnew:eq:minor}
\begin{pmatrix}
 1& 0& 0& 0& 0& 0& 0& 0\\
 0& 1& 0& 0& 0& 0& 0& 0\\
 0&-2& 1& 0& 0& 0& 0& 0\\
 0& 0& 0& 0& 0& 0& 0&-1\\
 0& 0& 0& 0& 0&-1& 0& 0\\
 0& 0& 0& 0& 0& 4&-1& 0\\
-1& 1& 0&-1& 0& 1& 0&-1\\
 0&-2& 1& 4& 1&-7&-1&11
\end{pmatrix}.
\end{equation}
Its determinant is $1$, proving independence of the eight classes.
The six-transformation span is consequently the eight-dimensional
kernel of $\Omega$. The full linear-letter span lies between that
space and the same kernel, proving the stronger assertion too.
\end{proof}

\subsection{Two exact quotient identities}

The rank statement also supplies concise constructive reductions.
Write $F_{a,b}(z)=\Li_{a,b}(z,1)$ and
$F_{a,b,c}(z)=\Li_{a,b,c}(z,1,1)$, and put
$u=1-z$, $v=z/(z-1)$.
Here $\equiv$ denotes equality modulo products of positive-weight
hyperlogarithms and endpoint constants, of total weight six.
The identities use the common analytic continuation from $0<z<1$.

\begin{proposition}\label{depthnew:prop:congruences}
One has
\begin{align}
 F_{4,1,1}(z)\equiv{}&-\Li_6(z)+F_{5,1}(z)
 +\Li_6(u)-F_{5,1}(u)-\Li_6(v),
 \label{depthnew:eq:411}\\
 2F_{3,2,1}(z)+F_{3,1,2}(z)\equiv{}&
 -\Li_6(z)+F_{5,1}(z)+F_{4,2}(z)\notag\\
 &-10\Li_6(u)+10F_{5,1}(u)+F_{4,2}(u)\notag\\
 &-\Li_6(v)-F_{4,2}(v).
 \label{depthnew:eq:321312}
\end{align}
\end{proposition}
\begin{proof}
Expand the nine rows obtained by applying $I,R,M$, in that order,
to $000001,000011,000101$.
The coefficient vectors for \eqref{depthnew:eq:411} and
\eqref{depthnew:eq:321312}, respectively, are
\[
 (-1,1,0,1,-1,0,-1,0,0),\qquad
 (-1,1,1,-10,10,1,-1,0,-1).
\]
Subtract each corresponding word combination from its left side.
Pairing the residual with each of the nine standard Lyndon Lie
polynomials gives zero, by direct expansion. The primitive space is dual
to the indecomposable quotient, so each residual belongs to $J^2$.
Landen pullback and the endpoint-transport formula then give the stated
analytic congruences. This is an equality in the specified quotient;
no product term is silently asserted to vanish as a function.
\end{proof}

The supplied standard-library verifier reconstructs all of these objects,
checks the determinant-one minor, verifies the two congruences, and
checks $\Omega$ against all $320$ nonempty binary shuffle products of
weight six. Every operation uses exact integers or rational numbers.
No numerical polylogarithm value or rank estimate enters.

\paragraph{Consequences and limits.}
The general complementary-depth bound three is sharp for this formal
class at weight six, even after the argument alphabet is enlarged by
all six standard transformations. Exactly one triple coordinate remains,
rather than an unspecified failure of a search.
At a particular cyclotomic point, additional identities can still kill
that coordinate. The theorem therefore proves neither numerical
irreducibility of $F_{3,2,1}(i)$ nor an obstruction to identities using
other color alphabets. The existing octahedral $S_4$ proof illustrates
why enlarging the available laws can overcome a restricted formal
obstruction. The manuscript's $S_6$ candidate remains conjectural.
